\begin{center}
	\large{\textbf{KATA PENGANTAR}}
\end{center}
\indent Alhamdulillahi Rabbil 'Alamin. Segala puja dan puji syukur dipanjatkan kehadirat Allah SWT karena telah melimpahkan rahmat, hidayah, petunjuk, kekuatan, dan rezeki kepada penulis. Tak lupa sholawat serta salam penulis haturkan kepada junjungan dan panutan mulia Nabi Muhammad SAW yang membawa dunia dari jalan kegelapan menuju jalan yang terang benderang. Dengan demikian, berkat hidayah Allah SWT serta risalah Nabi Muhammad SAW, penulis dapat menyelesaikan Tugas Akhir yang berjudul:
\begin{center}
	\section*{PERFORMA MESIN PANAS KUANTUM FRAKSIONAL BERBASIS PARTIKEL PADA KOTAK 1 DIMENSI}
\end{center}
sebagai salah satu syarat kelulusan Program Sarjana Fisika FSAD ITS.\\
\indent Penulisan Tugas Akhir ini dapat terselesaikan dengan baik berkat adanya bantuan dan bimbingan dari berbagai pihak. Penulis menyampaikan terima kasih kepada semua pihak yang telah membantu atas terselesaikannya Tugas Akhir ini.
\begin{enumerate}
    \item Tuhan yang maha Esa karena tanpa kasih sayang dan rahmat-Nya, penulis tidak akan bisa sampai di titik ini.
    \item Ayah, Ibu dan Adik yang senantiasa memberikan dukungan, doa, dan harapan kepada penulis.
    \item Bapak Dr. Heru Sukamto, M.Si. selaku dosen pembimbing dan dosen wali atas bimbingan, arahan, masukan, dan motivasinya kepada penulis.
    \item Bapak Prof. Drs. Agus Purwanto, M.Si., M.Sc., D.Sc. dan Bapak Malik Anjelh Baqiya, M.Si., Ph.D. selaku dosen penguji yang telah membantu penulis untuk melihat kekurangan dalam penulisan sehingga tugas akhir ini menjadi lebih baik.
    \item Bapak Dr.rer.nat Bintoro Anang Subagyo, M.Si. dan Bapak Dr. Lila Yuwana, M.Si. selaku dosen fisika teori yang telah memberikan begitu banyak ilmu, nasihat, dan motivasi.
    \item Bapak Dr. Gatut Yudhoyono, M.T selaku Kepala Departemen Fisika yang telah menjalankan departemen dengan baik sehingga proses kuliahan berjalan dengan lancar.
    \item Dosen-dosen departemen fisika secara umum yang telah memberikan beragam pengetahuan serta ilmu yang bermanfaat kepada penulis baik melalui kegiatan perkuliahan ataupun kegiatan non-perkuliahan.
    \item Teman-teman kolega seperjuangan Laboratorium Fisika Teori dan Filsafat Alam (LaFTiFA) yang senantiasa menjadi teman diskusi dan bercengkrama selama berada di Lab. Terima kasih atas waktu, ilmu, serta dukungannya.
    \item Seluruh pihak yang tidak bisa disebutkan satu-persatu yang selalu memberikan dukungan dan motivasi dalam menyelesaikan tugas akhir ini.
\end{enumerate}
Penulis menyadari bahwa dalam penulisan ini masih terdapat kekurangan, sehingga segala kritik dan saran yang sifatnya membangun sangat diharapkan untuk kesempurnaan Tugas Akhir ini. Penulis berharap  Tugas Akhir ini dapat bermanfaat bagi penulis sendiri pada khususnya dan pembaca pada umumnya.

\begin{flushright}
	Surabaya, Juli 2024\\
	\vspace{2cm}
	Penulis
\end{flushright}

% \newpage
% \thispagestyle{empty}
% \vspace*{\fill}
%     \begin{center}
% 	Halaman ini sengaja dikosongkan
%     \end{center}
% \vspace*{\fill}
% \pagebreak
%\begin{center}
%	Halaman ini sengaja dikosongkan
%\end{center}